\documentclass[10pt,letterpaper]{book} % , draft
\usepackage[utf8]{inputenc}
\usepackage[spanish, activeacute]{babel}
\usepackage{amsmath,amsfonts,amssymb,amsthm}
\usepackage[framed,numbered,autolinebreaks,useliterate]{mcode}
\usepackage{etex}
\usepackage[mode=buildnew]{standalone}
\usepackage{tikz}
    \usetikzlibrary{decorations.pathmorphing}
    \usetikzlibrary{arrows.meta}
    \usetikzlibrary{arrows}


\usepackage[left=4cm,top=2.5cm,right=2cm]{geometry}
% \usepackage{biblatex}
    \usepackage{anyfontsize}
%----------------------------------------Paquetes para graficos---------------------------------------%
%     \usepackage{graphicx} % para incluir graficos
    \usepackage{graphicx}
    \usepackage{float}          % para posiciones bien las figure con \begin{figure}[H]
    \usepackage{sidecap}        % caption al lado de la imagen
    \usepackage{wrapfig}        % figuras rodeadas de texto
    \usepackage{caption}        % caption para subimagenes
    \usepackage{subcaption}

%-----------------------------------------------------------------------------------------------------%

%----------------------------------------Paquetes varios ---------------------------------------------%
    \usepackage{draftwatermark}        % marca de agua
    \usepackage{color}              % permite escribir con colores.
    %\usepackage[usenames,dvipsnames,svgnames,table]{xcolor}    % alternativa para escribir con colores adicionales, ver http://en.wikibooks.org/wiki/LaTeX/Colors
    \usepackage{multicol}           % permite poner columnas en las enumeraciones.
    \usepackage{enumerate}          % permite cambiar al enumeracion asi: \begin{enumerate}[a)].
    \usepackage[bookmarks   = false,
        hidelinks,
%       colorlinks  = true,
%       linktoc     = black,
        linkcolor   = black]{hyperref}    % permite agregar hipervinculos a las referencias, pag web, correoes, etc. SOLO FUNCIONA EN PDF.
    \usepackage{lscape}         % para rotar cosas en 90° (landscape)
%-----------------------------------------------------------------------------------------------------%

%-------------------------------------- Opciones de las paginas --------------------------------------%
    % Dim hoja carta: 21.59 x 27.94 centimetros
% %     \addtolength{\topmargin}{-2.5cm}
% %     \addtolength{\evensidemargin}{-1.795cm}
% %     \addtolength{\oddsidemargin}{-1.795cm}
% %     \addtolength{\footskip}{0.4cm}
% %     \textheight=23cm                % Establece el largo del texto en cada página. El default es 19 cm.
% %     \textwidth=16cm                 % Establece el ancho del texto en cada página (en este caso, de 17 cm). El default es 14 cm.
    \setlength{\parskip}{5pt}           % Espacio entre parrafos
    \setcounter{secnumdepth}{3}
    \setcounter{tocdepth}{1}            % profundidad del arbol del indice
%-----------------------------------------------------------------------------------------------------%

%-------------------------------------- Opciones varias ----------------------------------------------%
%   \pagestyle{headings}                     % ???
%    \pagestyle{empty}                       % borrar numero de pag ???\documentclass[11pt,letterpaper]{article}
%    \boldmath                               % simbolos matematicos en negrita
%    \renewcommand{\baselinestretch}{1.5}    % interlinedo simple
%    \setlength{\parindent}{0pt}             % Sin sangria.
% --- marca de agua -------%
    \SetWatermarkFontSize{10pt}
    \SetWatermarkScale{10}
    \SetWatermarkLightness{0.9}     % intencidad del color
    \SetWatermarkText{\parbox{2.5cm}{%
  BORRADOR\vspace*{-8pt }\\\vspace*{-8pt }
  \href{mailto:neopazo@uc.cl}{\tiny neopazo@uc.cl}
 \\ 
  {\tiny \today}}}
%-----------------------------------------------------------------------------------------------------%

%-----------------------------------------------------------------------------------------------------%
%--------- Mis teoremas ---------%

\renewcommand{\qedsymbol}{$\blacksquare$}

\theoremstyle{plain}
    \newtheorem{prop}{Proposici\'on}[section]
    \newtheorem{teor}{Teorema}[section]
    \newtheorem{corol}{Corolario}[section]
    \newtheorem{lema}{Lema}[section]
    \newtheorem{obs}{Observaci\'on}[section]
\theoremstyle{definition}
    \newtheorem{defin}{Definici\'on}[section]
    \newtheorem{ejem}{Ejemplo}[section]
    \newtheorem*{nota}{Nota}
    \newtheorem*{notac}{Notaci\'on}
  
%--------------------------------%
%-------------------------------------- MIS DEFINICIONES -------------------------------------------
    \newtheorem{ejer}{Ejercicio}
\theoremstyle{remark}%
%-----------------------------------------------------------------------------------------------------%


\title{\includegraphics[width=3cm]{./img/escudo.eps}\\
{\small\bf  UNIVERSIDAD DE CONCEPCI\'ON}\\
{\small\bf    FACULTAD DE CIENCIAS F\'ISICAS Y MATEM\'ATICAS}\\
{\small\bf    DEPARTAMENTO DE INGENIER\'IA MATEM\'ATICA}\\
     \hskip 1pt \\
     \hskip 1pt \\
    {\bf\huge Title}
\hskip 1pt \\
     \hskip 1pt \\
    {\it Memoria para optar al t\'itulo de Ingeniero Civil Matem\'atico}\\
}
\author{Autor}

\date{{\bf 15 de septiembre de 2015}\\
    CONCEPCI\'ON, CHILE}

\DeclareFontShape{OT1}{cmr}{bx}{sc}{<-> cmbcsc10}{}
\begin{document}
\pagenumbering{roman}
% !TEX root = ../main.tex
% linea de arriba: hace que los  compiladores compilen el documento  main.tex al tratar de compilar este documento

\begin{titlepage}
\  \\ \vspace{0.5cm} 
\begin{center}
  {\includegraphics[width=3cm]{./img/escudo.pdf}
  %----------------------- ESCUDO --------------------------------------% 
  %----------------------- ESCUDO --------------------------------------%
  \\ \vspace{0.4cm} 
  {\small\bf  UNIVERSIDAD DE CONCEPCI\'ON}\\ \vspace{0.4cm}
  {\small\bf    FACULTAD DE CIENCIAS F\'ISICAS Y MATEM\'ATICAS}\\ \vspace{0.4cm}
  {\small\bf    DEPARTAMENTO DE INGENIER\'IA MATEM\'ATICA}\\ \vspace{0.4cm}
      \hskip 1pt \\ \vspace{0.4cm} 
      \hskip 1pt \\ \vspace{0.4cm}
	  {\bf\huge Titulo}
  \hskip 1pt \\ \vspace{0.8cm}
      \hskip 1pt \\ %\vspace{1cm}
      {\it\LARGE Memoria para optar al t\'itulo de Ingeniero Civil Matem\'atico}\\ \vspace{0.5cm}}
  {
  \large
  \begin{center}
  \begin{tabular}{rl}
  \vspace{0.5cm} 
  Autor:			& {\bf Autor}\\
  Profesores Guía:	& {\bf Prof 1} \\ 			
					& {\bf Prof 2}
  \end{tabular}
  \end{center}
  
  %   {Autor: {\bf Diego Alberto Maldonado Montiel}\\ \vspace{0.5cm}
  %   Profesores Guía: {\bf Anahí Gajardo Schulz}  \\ \vspace{0.5cm} \hspace{3.5cm}{\bf Andrés Moreira Wenzel}
  \vfill
  {{\bf 15 de septiembre de 2015}\\ 
  CONCEPCI\'ON, CHILE}
  }
\end{center}
\end{titlepage} 

   \renewcommand{\baselinestretch}{1}    % interlinedo simple

\cleardoublepage

\vspace{4cm}
\setcounter{page}{3}
\begin{center}
\large{\textbf{Agradecimientos}} 
\end{center}


{\it 
Quiero agradecer a ...

También quiero agradecer a ...

Asimismo, agradezco a ...

Finalmente quiero agradecer el apoyo económico de ... 


}

\tableofcontents % indice de contenidos

\cleardoublepage
\addcontentsline{toc}{chapter}{Lista de figuras} % para que aparezca en el indice de contenidos
\listoffigures % indice de figuras

\cleardoublepage
\addcontentsline{toc}{chapter}{Lista de tablas} % para que aparezca en el indice de contenidos
\listoftables % indice de tablas

\thispagestyle{empty}
\cleardoublepage
\pagenumbering{arabic} 
\chapter{Introducci\'on}
 % !TEX root = ../main.tex
% linea de arriba: hace que los  compiladores compilen el documento  main.tex al tratar de compilar este documento

Introducimos ...

\chapter{Preliminares}
 \input{./partes/preliminares.tex}
\chapter{Desarrollo}
 % !TEX root = ../main.tex
% linea de arriba: hace que los  compiladores compilen el documento  main.tex al tratar de compilar este documento
Desarrollo \cite{wikipedia1}

\begin{figure}[H]
  \centering
  \begin{subfigure}[b]{0.475\textwidth}
    \centering
        {
          \includegraphics[width= 0.9\textwidth]{img/figura.pdf} % pdf
        }
        \caption{$t=1$}     
  \end{subfigure}
  \begin{subfigure}[b]{0.475\textwidth}
    \centering
        {
          \includestandalone[width=0.9\textwidth]{img/figura2} % tikz
        }     
        \caption{$t=2$}
  \end{subfigure}
\caption{Ejemplo de figura.}\label{fig: ejem}
\end{figure}

\begin{lstlisting}
for i = 1:3
  if i >= 5 && a ~= b       % literate programming replacement
    disp('cool');           % comment with some §\mcommentfont\LaTeX in it: $\mcommentfont\pi x^2$§
  end
  [:,ind] = max(vec);
  x_last = x(1,end) - 1;
  v(end);
  really really long really really long really really long really really long really really long line % blaaaaaaaa
  ylabel('Voltage (µV)');
end
\end{lstlisting}
\chapter{Conclusiones y trabajos futuros}
 % !TEX root = ../main.tex
% linea de arriba: hace que los  compiladores compilen el documento  main.tex al tratar de compilar este documento
\section{Conclusiones}


En estas conclusiones concluimos conclusiones como conclusion de nuestro trabajo. 
De este,  tambien concluimos otras conclusiones que pueden ser concluidas conmo conclusion de nuestras conlcusiones y que era esperado que fueran concluidas al concluir el trabajo.

Finalmente dejamos sin concluir algunas conclusiones, pues estas no pudieron ser concluidas por la investigacion realizada, debiendo ser concluidas por otros investigadores, quienes las conclurian en sus concluisones de la misma forma que nosotros concluimos  nuestras conlusiones, concluyendo el tema.
% \chapter{Anexo}
%  \input{./partes/anexo.tex}

\bibliographystyle{plainnat}
\bibliography{biblio}



 

\end{document}

    